\documentclass[../../../include/open-logic-section]{subfiles}

\begin{document}

\olfileid{sth}{cardinals}{classing}
\olsection{有限、可数与不可数基数}

在引入了基数之后，我们不妨花点时间来讨论不同种类的基数；
具体地说，即有限、可数与不可数基数。

我们的前两个结果蕴涵：有限基数恰好就是有限序数，
我们在 \olref[z][infinity-again]{defnomega} 中将后者定义为\emph{自然数}：

\begin{prop}\ollabel{finitecardisoequal}
令 $n, m \in \omega$。则 $n = m$ 当且仅当 $\cardeq{n}{m}$。
\end{prop}

\begin{proof}
\emph{从左到右}是平凡的。为证明\emph{从右到左}，
假设 $\cardeq{n}{m}$ 但 $n \neq m$。
由三歧律，要么 $n \in m$，要么 $m \in n$；不失一般性，设 $n \in m$。
那么 $n \subsetneq m$，并且存在一个双射 $f \colon m \to n$，
因此 $m$ 是 Dedekind（戴德金）无限的，
这与 \olref[z][infinity-again]{naturalnumbersarentinfinite} 矛盾。
\end{proof}

\begin{cor}\ollabel{naturalsarecardinals}
若 $n \in \omega$，则 $n$ 是基数。
\end{cor}

\begin{proof}
直接可得。
\end{proof}
\noindent
由此还可得出，描述基数何为“有限”或“无限”的几种合理概念是一致的：
\begin{thm}\ollabel{generalinfinitycharacter}对任意集合 $A$，以下陈述等价：
\begin{enumerate}
	\item\ollabel{card:notinomega} $\card{A} \notin \omega$，即
	$A$ 不是一个自然数；
	\item\ollabel{card:omegaplus} $\omega \leq \card{A}$；
	\item\ollabel{card:infinite} $A$ 是 Dedekind 无限的。
\end{enumerate}
\end{thm}

\begin{proof}
由 \olref[ord-arithmetic][using-addition]{ordinfinitycharacter}、
\olref[cardinals][cardsasords]{lem:CardinalsBehaveRight} 和
\olref{naturalsarecardinals} 得出。
\end{proof}

据此，我们可以为以下在 \olref[sfr][][]{part} 中颇为非正式地使用过的概念，给出如下\emph{定义}：

\begin{defn}\ollabel{defnfinite}
我们称 $A$ 是\emph{有限的}，当且仅当 $\card{A}$ 是一个自然数，即 $\card{A} \in \omega$。
否则，我们称 $A$ 是\emph{无限的}。
\end{defn}
\noindent 
但需注意，此定义是在 $\ZFC$ 的背景下给出的。毕竟，我们需要良序定理来保证每个集合都有基数。
事实上，没有良序定理，就可能存在既非有限亦非 Dedekind 无限的集合。
我们将在 \olref[choice][]{chap} 中回到这类问题。目前，我们继续依赖良序定理。

现在，让我们从有限基数转向无限基数。
以下是两个基本结论：

\begin{cor}\ollabel{omegaisacardinal}
$\omega$ 是最小的无限基数。
\end{cor}

\begin{proof}
$\omega$ 是基数，因为 $\omega$ 是 Dedekind 无限的，且若对任意 $n \in \omega$ 有 $\cardeq{\omega}{n}$，
则 $n$ 将是 Dedekind 无限的，
与 \olref[z][infinity-again]{naturalnumbersarentinfinite} 矛盾。
由定义，$\omega$ 就是最小的无限基数。
\end{proof}

\begin{cor}
每个无限基数都是一个极限序数。
\end{cor}

\begin{proof}
设 $\alpha$ 是一个无限后继序数，则存在某个 $\beta$ 使得 $\alpha = \beta
\ordplus 1$。由 \olref{finitecardisoequal}，$\beta$
也是无限的，所以根据 \olref[ord-arithmetic][using-addition]{ordinfinitycharacter} 有 $\cardeq{\beta}{\beta \ordplus 1}$。
由 \olref[cardinals][cardsasords]{lem:CardinalsBehaveRight} 可得 $\card{\beta} = \card{\beta\ordplus 1} = \card{\alpha}$，
因此 $\alpha \neq \card{\alpha}$。
\end{proof}

早在 \olref[sfr][siz][enm-alt]{defn:enumerable} 中，我们就曾指出可以区分可数与不可数的无限集合。该定义自然引出以下结论：

\begin{prop}
$A$ 是可数的当且仅当 $\card{A} \leq \omega$，而 $A$ 是不可数的当且仅当 $\omega < \card{A}$。
\end{prop}

\begin{proof}
由三歧律，这两个断言等价，因此只需证明 $A$ 可数当且仅当 $\card{A} \leq \omega$。
\emph{从右到左}：若 $\card{A} \leq \omega$，则由
\olref[cardinals][cardsasords]{lem:CardinalsBehaveRight} 和
\olref{omegaisacardinal} 可得 $\cardle{A}{\omega}$。
\emph{从左到右}：假设 $A$ 可数；则根据 \olref[sfr][siz][enm-alt]{defn:enumerable}，有三种可能情况：
\begin{enumerate}
	\item 若 $A = \emptyset$，则根据 \olref{naturalsarecardinals} 和
	\olref[cardinals][cardsasords]{lem:CardinalsBehaveRight}，有 $\card{A} = 0 \in \omega$。
	\item 若 $\cardeq{n}{A}$，则根据 \olref{naturalsarecardinals} 和
	\olref[cardinals][cardsasords]{lem:CardinalsBehaveRight}，有 $\card{A} = n \in \omega$。
	\item 若 $\cardeq{\omega}{A}$，则根据 \olref{omegaisacardinal}，有 $\card{A} = \omega$。
\end{enumerate}
因此在所有情况下，都有 $\card{A} \leq \omega$。 
\end{proof}
\noindent
事实上，$\omega$ 具有特殊地位。尽管存在许多可数序数：

\begin{cor}
$\omega$ 是唯一的可数无限基数。
\end{cor}

\begin{proof}
设 $\cardfont{a}$ 是一个可数无限基数。由于 $\cardfont{a}$ 是无限的，有 $\omega \leq \cardfont{a}$。由于 $\cardfont{a}$ 是可数基数，$\cardfont{a} =
\card{\cardfont{a}} \leq \omega$。故由三歧律，$\cardfont{a} = \omega$。
\end{proof}

当然，存在无限多个基数。于是你或许会问：\emph{有多少个基数呢？} 以下结果表明，我们可能需要重新考虑这个问题。

\begin{prop}\ollabel{unioncardinalscardinal}
如果 $X$ 的每个成员都是基数，那么 $\bigcup X$ 也是基数。
\end{prop}

\begin{proof}
容易验证 $\bigcup X$ 是一个序数。设 $\alpha \in \bigcup X$ 是一个序数；则存在某个基数 $\cardfont{b}$，使得 $\alpha \in \cardfont{b} \in X$。由于 $\cardfont{b}$ 是基数，有 $\cardless{\alpha}{\cardfont{b}}$。由于 $\cardfont{b} \subseteq \bigcup X$，我们有 $\cardle{\cardfont{b}}{\bigcup X}$，因此 $\cardneq{\alpha}{\bigcup X}$。推广可知，$\bigcup X$ 是一个基数。
\end{proof} 

\begin{thm}\ollabel{lem:NoLargestCardinal}
不存在最大的基数。
\end{thm}

\begin{proof}
对任意基数 $\cardfont{a}$，Cantor（康托尔）定理 (\olref[sfr][siz][car]{thm:cantor}) 与 \olref[cardinals][cardsasords]{lem:CardinalsExist} 蕴涵了 $\cardfont{a} < \card{\Pow{\cardfont{a}}}$。
\end{proof}

\begin{thm}
所有基数构成的集合不存在。
\end{thm}

\begin{proof}
用归谬法，假设 $C = \Setabs{\cardfont{a}}{\cardfont{a} \text{
is a cardinal}}$。根据 \olref{unioncardinalscardinal}，$\bigcup C$ 是一个基数，因此由 \olref{lem:NoLargestCardinal} 可知，存在一个基数 $\cardfont{b} > \bigcup C$。根据定义，$\cardfont{b} \in C$，所以 $\cardfont{b} \subseteq \bigcup{C}$，从而 $\cardfont{b} \leq \bigcup C$，得出矛盾。
\end{proof}

应将此结论与 Russell（罗素）悖论及 Burali-Forti（布拉利-福尔蒂）悖论进行比较。

\end{document}